\section{Discussion, limitations, and future work}
\label{sec:discussion}

\subsection{What the split buys}

The CPU octree is a semantic store: it knows the SDF, the brush, the
mapping, the version, and the LoD ladder.  The GPU mesher is a pure
function from a flattened neighbourhood to a dual mesh.  This separation
makes both sides simpler:
\begin{itemize}
\item the CPU never needs a triangle mesh to answer ``what is the
  material at this cell?'';
\item the GPU never needs a pointer, an allocator, or a tree lock;
\item the renderer never needs a conversion pass;
\item the exporter can reuse the CPU walker as a reference.
\end{itemize}

\subsection{The remaining host cost}

The host still pays for the flattened upload and the two fence waits per
chunk.  The measured 8--19 ms/chunk is not negligible.  The cost is on
worker threads, so it does not directly stall the render thread, but it
consumes CPU cores and memory bandwidth.  The next step would be to batch
several chunks into one submission and share one descriptor set, as the
perf reports propose.  That would amortize the submit/fence cost and the
descriptor writes.  The current per-chunk design was chosen for
incremental streaming and easy failure isolation; batching is a
throughput optimization, not a correctness requirement.

\subsection{The mesh is not the only consumer}

The same span feeds:
\begin{itemize}
\item the raster solid pass;
\item the shadow cascades;
\item the ray-tracing BLAS;
\item the vegetation placement compute;
\item the SDF fire anchors (via node SDFs).
\end{itemize}
This is why the span protocol is centralized in \code{IndirectRenderer}.
If each consumer allocated its own copy, the memory and the synchronization
surface would explode.  The packed pool is the single source of geometry
truth for the frame.

\subsection{Adaptive LoD and cracks}

The adaptive stitching is exact for the chunk-local walk, but it relies on
both chunks seeing the same four quadrant cells.  The one-cell margin
guarantees this for a neighbour that is at most one cell larger.  A
pathological tree with a much larger neighbour (more than one cell) could
still sample outside the margin.  The flattening uses the largest target
cell in the subtree, not the chunk's own cell, precisely to cover this
case.  A future change that shrinks the margin must re-run the
CPU-vs-GPU self-check on a chunk with a coarse neighbour.

\subsection{The coarse-rung density cap}

Vegetation placement caps a single coarse rung at a fixed instance count.
This is a deliberate visual/performance trade: a root-scale chunk can
request hundreds of thousands of tufts, which is both expensive and
visually indistinguishable from a smaller count at its distance.  The cap
is logged, not silent.  A future distance-density model could replace the
hard cap with a screen-space or distance-based thinning curve.

\subsection{The fire anchors}

The fire effect now derives anchors from the node SDFs rather than from the
triangulated surface.  This is cheaper and fits the GPU-only mesh, but it
is coarser than the old triangle sampling.  The Newton step lands on the
trilinear zero crossing; for a primitive sphere or cylinder this is
accurate.  For a highly non-planar cell, the anchor can be a few
centimetres off the true QEF surface.  The flame scale is large enough
that this is not visible.  A future improvement could run a small GPU
placement pass over the brush-4 triangles and read back only the anchor
records.

\subsection{What would break the design}

The following changes would violate the contract and require a redesign:
\begin{itemize}
\item reading a vertex back to the host (destroys the no-readback
  property);
\item letting the CPU write a chunk vertex (destroys the single-producer
  invariant);
\item freeing a span without a cull-frame delay (use-after-free);
\item changing the node table order without re-running the exporter
  comparison (changes slot assignment and the BLAS);
\item dropping a triangle on overflow instead of failing (silent holes).
\end{itemize}

\subsection{Future work}

Three directions are natural:
\begin{enumerate}
\item \textbf{Batched topology.}  Flatten all chunks of a layer into one
  arena, record one geometry dispatch for the batch, and read all counters
  with one map.  This is the highest-leverage performance change.
\item \textbf{GPU-side fire placement.}  A compute pass over brush-4
  triangles could produce anchors in a device buffer and the SDF scene
  could consume them directly, avoiding the host anchor list.
\item \textbf{Compressed vertex payload.}  The 64-byte vertex is generous;
  a packed normal/brush/mapping format could halve the pool and the
  bandwidth, at the cost of a decode in the vertex shader and the BLAS
  build.
\end{enumerate}
